% ============================================================================
%  CAPÍTULO II: OBJETIVOS DEL TRABAJO
% ============================================================================

\section{Objetivos del trabajo}

\subsection{Objetivo general}

Comprender el funcionamiento de un servidor de correo electrónico (SMTP/IMAP/POP3)
y un servidor de archivos (SMB/CIFS), su configuración e integración con los
servicios de resolución de nombres (DNS) y autenticación de usuarios en un
entorno de red administrado con Debian.

\subsection{Objetivos específicos}

\begin{itemize}
  \item Comprender la arquitectura y el funcionamiento del protocolo SMTP,
        incluyendo los roles de MTA y la transferencia de mensajes entre
        servidores.
  \item Implementar un servidor de correo electrónico funcional utilizando
        Postfix (envío) y Dovecot (recepción) instalados directamente
        sobre Debian.
  \item Configurar los registros DNS (MX y A) necesarios para la resolución
        del dominio de correo dentro de la red interna.
  \item Configurar la autenticación de usuarios mediante SASL y la protección
        de las conexiones con STARTTLS mediante certificados auto-firmados.
  \item Instalar y configurar Samba (SMB) para compartir carpetas en la red
        local, asignando permisos de acceso por usuario.
  \item Documentar el proceso completo de instalación, configuración y prueba
        de ambos servicios, incluyendo comandos, archivos de configuración y
        evidencias de cada paso.
\end{itemize}
